\subsection{Event Code and Sources}\label{section:event-code-and-sources}
Every delivered event carries a 32-bit code. The high byte is the event class and the low 24 bits are an ID in that
class\textquotesingle{}s namespace.

(:event-classes:base.events.EXCEPTION,base.events.INTERRUPT,base.events.NMI:)

The following fixed event codes are assigned by this specification. Each row answers only which symbolic event is
assigned to a fixed code; the event's behavior, frame, and (:term:base.terms.payload:) are defined by the topic that owns that behavior.

\BedrockTableCaption{Fixed Architectural Event-Code Assignments}
\begin{BedrockLongTable}{@{}>{\raggedright\arraybackslash}p{1.10in}>{\raggedright\arraybackslash}p{4.55in}@{}}
\toprule
\rowcolor{BedrockHeaderFill}
\textbf{Code} & \textbf{Event}\\
\midrule
\endfirsthead
\multicolumn{2}{@{}l}{\scriptsize\itshape Table \theBedrockTable\ (continued)}\\
\toprule
\rowcolor{BedrockHeaderFill}
\textbf{Code} & \textbf{Event}\\
\midrule
\endhead
(:event-code:base.events.EXCEPTION.SINGLE_STEP:)
(:event-code:base.events.EXCEPTION.EXPLICIT_BREAKPOINT:)
(:event-code:base.events.EXCEPTION.EXECUTION_BREAKPOINT:)
(:event-code:base.events.EXCEPTION.MEMORY_READ_WATCHPOINT:)
(:event-code:base.events.EXCEPTION.MEMORY_WRITE_WATCHPOINT:)
(:event-code:base.events.EXCEPTION.ATOMIC_MEMORY_WATCHPOINT:)
(:event-code:base.events.EXCEPTION.PRIVILEGE_VIOLATION:)
(:event-code:base.events.EXCEPTION.SYSTEM_CALL:)
(:event-code:base.events.EXCEPTION.DIVIDE_BY_ZERO:)
(:event-code:base.events.EXCEPTION.SIGNED_DIVIDE_OVERFLOW:)
(:event-code:base.events.EXCEPTION.INVALID_OPCODE:)
(:event-code:base.events.EXCEPTION.INVALID_ADDRESSING_FORM:)
(:event-code:base.events.EXCEPTION.RESERVED_INSTRUCTION_ENCODING:)
(:event-code:base.events.EXCEPTION.UNAVAILABLE_INSTRUCTION_EXTENSION:)
(:event-code:base.events.EXCEPTION.EXPLICIT_ILLEGAL_INSTRUCTION:)
(:event-code:base.events.EXCEPTION.TRUNCATED_INSTRUCTION:)
(:event-code:base.events.EXCEPTION.INVALID_OPERAND_RELATION:)
(:event-code:base.events.EXCEPTION.PAGE_NOT_PRESENT:)
(:event-code:base.events.EXCEPTION.PAGE_PERMISSION_VIOLATION:)
(:event-code:base.events.EXCEPTION.MALFORMED_PAGE_TABLE_ENTRY:)
(:event-code:base.events.EXCEPTION.NONCANONICAL_ADDRESS:)
(:event-code:base.events.EXCEPTION.SEGMENT_BOUNDS_VIOLATION:)
(:event-code:base.events.EXCEPTION.ATOMIC_ALIGNMENT_FAULT:)
(:event-code:base.events.EXCEPTION.MEMORY_TYPE_FAULT:)
(:event-code:base.events.EXCEPTION.PHYSICAL_ADDRESS_FAULT:)
(:event-code:base.events.EXCEPTION.MMIO_ALIGNMENT_FAULT:)
(:event-code:base.events.EXCEPTION.UNSUPPORTED_MMIO_OPERATION:)
(:event-code:base.events.EXCEPTION.INVALID_CONTROL_SELECTOR:)
(:event-code:base.events.EXCEPTION.RESERVED_CONTROL_BITS:)
(:event-code:base.events.EXCEPTION.INVALID_CONTROL_IMAGE:)
(:event-code:base.events.EXCEPTION.INVALID_CONTROL_TRANSITION:)
(:event-code:CFI.events.EXCEPTION.CONTROL_FLOW_INTEGRITY_VIOLATION:)
(:event-code:VECTOR.events.EXCEPTION.VECTOR_LOOP_OFFSET_OUT_OF_RANGE:)
(:event-code:VECTOR.events.EXCEPTION.VECTOR_LANE_INDEX_OUT_OF_RANGE:)
(:event-code:FP.events.EXCEPTION.FLOATING_POINT_EXCEPTION:)
(:event-code:base.events.EXCEPTION.BUS_NO_RESPONDER:)
(:event-code:base.events.EXCEPTION.BUS_ACCESS_DENIED:)
(:event-code:base.events.EXCEPTION.BUS_TIMEOUT:)
(:event-code:base.events.EXCEPTION.BUS_DATA_ERROR:)
(:event-code:base.events.EXCEPTION.BUS_OTHER_ERROR:)
(:event-code:base.events.EXCEPTION.EVENT_ENTRY_STATE_FAILURE:)
(:event-code:base.events.EXCEPTION.EVENT_STACK_STATE_FAILURE:)
(:event-code:base.events.EXCEPTION.EVENT_FRAME_ADDRESS_FAILURE:)
(:event-code:base.events.EXCEPTION.EVENT_FRAME_STORE_FAILURE:)
(:event-code:base.events.EXCEPTION.MACHINE_CHECK:)
\bottomrule
\end{BedrockLongTable}

Unassigned base-profile exception IDs are (:term:base.terms.reserved:). Interrupt identities do not share a namespace
with exception IDs. Interrupt-controller routing, native-ID mapping, ownership, trigger mode, and any external
acknowledgement or end-of-interrupt protocol belong to the platform ABI. A valid posted identity is recorded in the
target logical processor's architectural interrupt file; the delivered event code has class \texttt{0x01} and that
identity in its low 24 bits.

Deadline-timer expiry uses the same posting operation with the identity supplied by TARM. It does not define a
fixed timer identity or a separate interrupt-source catalog member.

A restartable fault saves the faulting instruction or its architecturally defined restart point. A trap saves the
following instruction. An interrupt or NMI saves the next instruction that would otherwise execute. A faulting repeat
iteration rolls back without decrementing its counter, then saves the repeat-prefix address. After a successful repeat
iteration, the counter decrement commits before an event may be admitted between iterations; that event also saves
the repeat-prefix address. Event entry clears hidden repeat state, and ERET executes the saved prefix normally without
restoring repeat (:term:base.terms.continuation:).
